% SPDX-License-Identifier: CC-BY-4.0
\documentclass[11pt,letterpaper]{article}

\usepackage[T1]{fontenc}
\usepackage{amsmath,amssymb}
\usepackage{tikz}
\usepackage{pgfplots}
\pgfplotsset{compat=1.18}
\usepackage{algorithm}
\usepackage{algpseudocode}
\usepackage{booktabs}
\usepackage{enumitem}
\usepackage[hidelinks]{hyperref}

\newcommand{\N}{\mathbb{N}}
\newcommand{\Z}{\mathbb{Z}}
\newcommand{\Q}{\mathbb{Q}}
\newcommand{\R}{\mathbb{R}}
\newcommand{\C}{\mathbb{C}}
\newcommand{\Map}[3]{#1 \colon #2 \to #3}
\newcommand{\Mod}[1]{#1\text{-}\mathrm{Mod}}
\newcommand{\GroupHom}{\mathrm{Hom}_{\mathbf{Grp}}}
\newcommand{\RingHom}{\mathrm{Hom}_{\mathbf{Ring}}}

\newtheorem{theorem}{Theorem}
\newtheorem{lemma}[theorem]{Lemma}
\newtheorem{corollary}[theorem]{Corollary}
\newtheorem{definition}[theorem]{Definition}
\newtheorem{remark}[theorem]{Remark}
\newenvironment{proof}{\noindent\textit{Proof.}}{\hfill\(\square\)}

\title{Generalized Collatz Library}
\author{}
\date{\today}

\begin{document}
\maketitle
\section*{Magnitude of one affine return}

\begin{abstract}
The absolute value of the affine state \(q^{j}n+c\,U_{j}(q)\) is at most
\(|q|^{j}|n|+|c|\,|U_{j}(q)|\). The only inequality is the triangle inequality.
The absolute value of a product and the absolute value of a power are
equalities, so they preserve the direction. If that state is the first affine
state divisible by the signed divisor \(d\), exact division gives
\(|d|\,|(q^{j}n+c\,U_{j}(q))/d|\) equal to the absolute value of the state, and
the same majorant bounds the quotient. The comparison is one seed and one
hitting time. It is not a uniform core.
\end{abstract}

\section{Affine states}

Throughout, \(q\), \(c\), and \(n\) denote integers and \(j\) denotes a natural
number, including zero. The offset sum is defined by \(U_{0}(q)=0\) and
\[
U_{j+1}(q)=U_{j}(q)+q^{j}.
\]
Thus \(U_{j}(q)\) is the sum of \(q^{0},\ldots,q^{j-1}\) when \(j\ge 1\), and
the empty sum is zero. No closed form of the sum is required.

\begin{definition}
\label{def:affine-state}
The \emph{affine state} at time \(j\) is
\[
A_{j}(q,c,n)=q^{j}n+c\,U_{j}(q).
\]
\end{definition}

A parameter record consists of integers \(q\), \(c\), and \(d\) with
\(|d|\ge 2\). Then \(d\neq 0\), since \(|d|\ge 2>0\). The raw map sends an
integer \(m\) to \(m/d\) when \(d\) divides \(m\), and to \(qm+c\) otherwise.
The theorems below do not iterate that map past a first hit. They bound the
integer \(A_{j}(q,c,n)\) at one time at which \(d\) divides it.

\begin{definition}
\label{def:first-hit}
Let the parameter record have multiplier \(q\), offset \(c\), and divisor
\(d\). The time \(j\) is a \emph{first hit} of the seed \(n\) when
\(d\) divides \(A_{j}(q,c,n)\) and \(d\) divides \(A_{i}(q,c,n)\) for no
natural number \(i<j\).
\end{definition}

When \(j=0\), there is no natural number \(i<0\), so a first hit at time zero
is divisibility of the seed. The definition does not assert that a first hit
exists.

\section{The triangle bound}

\begin{theorem}
\label{thm:affine-triangle}
For all \(q,c,n\in\Z\) and every \(j\in\N\),
\[
\bigl|q^{j}n+c\,U_{j}(q)\bigr|
\le |q|^{j}\,|n|+|c|\,\bigl|U_{j}(q)\bigr|.
\]
\end{theorem}

\begin{proof}
Write \(a=q^{j}n\) and \(b=c\,U_{j}(q)\). The triangle inequality on \(\Z\)
states that \(|a+b|\le |a|+|b|\). This is the only step that is an inequality.

The absolute value of a product is the product of the absolute values, so
\[
|a|=|q^{j}n|=|q^{j}|\,|n|
\]
and
\[
|b|=|c\,U_{j}(q)|=|c|\,\bigl|U_{j}(q)\bigr|.
\]
For every integer \(x\) and every natural number \(k\), including zero,
\(|x^{k}|=|x|^{k}\). Indeed, \(x^{0}=1\) and \(|x|^{0}=1\). If the identity
holds at \(k\), then
\[
|x^{k+1}|=|x^{k}x|=|x^{k}|\,|x|=|x|^{k}\,|x|=|x|^{k+1}.
\]
Applied to \(x=q\) and \(k=j\), this gives \(|q^{j}|=|q|^{j}\). Therefore
\[
|a|+|b|=|q|^{j}\,|n|+|c|\,\bigl|U_{j}(q)\bigr|.
\]
Substituting this equality into the right-hand side of
\(|a+b|\le |a|+|b|\) does not change that side. The resulting comparison is
still \(\le\), not \(\ge\). The left-hand side is
\(\bigl|q^{j}n+c\,U_{j}(q)\bigr|\).
\end{proof}

The two terms on the left may cancel. The inequality may therefore be strict.
Equality holds when \(q^{j}n\) and \(c\,U_{j}(q)\) are both nonnegative or both
nonpositive, but the theorem does not require that alignment.

\section{Exact division}

\begin{theorem}
\label{thm:exact-division}
Let \(d,a\in\Z\), and assume \(d\) divides \(a\). Let \(a/d\) denote the
integer quotient satisfying \(d\cdot(a/d)=a\). Then
\[
|d|\,\bigl|a/d\bigr|=|a|.
\]
\end{theorem}

\begin{proof}
Divisibility supplies an integer witness \(m\) with \(a=dm\). The quotient in
the statement is that witness: \(m=a/d\), and the displayed product identity
is \(d\cdot(a/d)=a\). The absolute value of a product then gives
\[
|a|=|d\cdot(a/d)|=|d|\,\bigl|a/d\bigr|.
\]
Both steps after the choice of quotient are equalities. If \(d\neq 0\), the
witness \(m\) is unique. A parameter divisor satisfies \(|d|\ge 2\), hence
\(d\neq 0\), so the quotient in the return bound below is that unique integer.
The identity itself uses only \(d\cdot(a/d)=a\), not a lower bound on \(|d|\).
\end{proof}

\section{One return}

\begin{theorem}
\label{thm:return-magnitude}
Let \(q\), \(c\), and \(d\) be a parameter record, so \(|d|\ge 2\). Let
\(n\in\Z\) and \(j\in\N\). If \(j\) is a first hit of the seed \(n\), then
\[
|d|\,\Bigl|\frac{A_{j}(q,c,n)}{d}\Bigr|
\le |q|^{j}\,|n|+|c|\,\bigl|U_{j}(q)\bigr|.
\]
The left-hand side equals \(\bigl|A_{j}(q,c,n)\bigr|\) by
Theorem~\ref{thm:exact-division}. The right-hand side is the majorant of
Theorem~\ref{thm:affine-triangle}. The statement concerns one seed and one
hitting time. It is not a uniform core.
\end{theorem}

\begin{proof}
Write \(A=A_{j}(q,c,n)\). By Definition~\ref{def:first-hit}, \(d\) divides
\(A\). Theorem~\ref{thm:exact-division}, applied to this divisor and this
integer, yields
\[
|d|\,\bigl|A/d\bigr|=|A|.
\]
Definition~\ref{def:affine-state} expands the right-hand absolute value as
\[
|A|=\bigl|q^{j}n+c\,U_{j}(q)\bigr|.
\]
Theorem~\ref{thm:affine-triangle}, at the same \(q\), \(c\), \(n\), and \(j\),
yields
\[
|A|\le |q|^{j}\,|n|+|c|\,\bigl|U_{j}(q)\bigr|.
\]
The preceding display is an equality, so replacing \(|A|\) by
\(|d|\,\bigl|A/d\bigr|\) preserves the inequality direction. This is the
required comparison.

The first-hit hypothesis also says that \(d\) divides \(A_{i}(q,c,n)\) for no
\(i<j\). That clause is not an input to the estimate. The estimate uses
divisibility at the single time \(j\), the exact-quotient identity at that
time, and the triangle inequality at that same time. Retaining the clause
keeps the hypothesis equal to a first hit rather than to divisibility at an
arbitrary time.

No second seed appears. No second time appears. The majorant still depends on
\(q\), \(j\), \(n\), \(c\), and \(U_{j}(q)\). The argument does not take a
supremum over seeds, over hitting times, or over a family of blocks, and it
does not produce a bound independent of the seed and the time. In that sense
the comparison is not a uniform core.
\end{proof}

\begin{remark}
The estimate uses divisibility only at time \(j\). The clause that no earlier
affine state is divisible is retained so that the hypothesis is a first hit,
and it is not used as an input to either absolute value. The quantifiers remain
one parameter record, one seed, and one time. Nothing is asserted about a raw
state after the hit, about existence of a hit, or about a bound on the hitting
time.
\end{remark}

\end{document}
